\documentclass{article}

\usepackage[utf8]{inputenc}
\usepackage[T1]{fontenc}
\usepackage{amsmath,amssymb,amsthm}
\usepackage{mathtools}
\usepackage{hyperref}
\usepackage{booktabs}

% Lightweight citation command for this standalone draft (no natbib).
% Renders \citep{key} as (Paper N) based on the key.
\makeatletter
\newcommand{\citep}[2][]{%
  \begingroup
  \def\@tmpkey{#2}%
  \ifx\@tmpkey\@empty\else
    (\ifx\@empty#1\@empty\else #1, \fi
     \@nameuse{bibkey@\@tmpkey})%
  \fi
  \endgroup
}
\expandafter\def\csname bibkey@lasser2026gfm\endcsname{Paper~1}
\expandafter\def\csname bibkey@lasser2026poset\endcsname{Paper~2}
\expandafter\def\csname bibkey@lasser2026horizon\endcsname{Paper~3}
\expandafter\def\csname bibkey@lasser2026scm\endcsname{Paper~4}
\expandafter\def\csname bibkey@lasser2026exo\endcsname{Paper~5}
\expandafter\def\csname bibkey@lasser2026wamura\endcsname{Paper~6}
\expandafter\def\csname bibkey@lasser2026phase\endcsname{Paper~7}
\makeatother

\newtheorem{definition}{Definition}
\newtheorem{proposition}{Proposition}
\newtheorem{theorem}{Theorem}
\newtheorem{conjecture}{Conjecture}
\newtheorem{remark}{Remark}
\newtheorem{example}{Example}

\DeclareMathOperator{\vol}{vol}
\newcommand{\volL}{{\vol_{\mathrm{P}}}}
\newcommand{\volR}{{\vol_{\mathrm{R}}}}
\newcommand{\volRlower}{{\vol_{\mathrm{R}}^{\mathrm{lower}}}}
\newcommand{\canon}{\mathrm{canon}}
\newcommand{\ResS}{R_{\mathrm{S}}}

\title{Draft v4.1: Capability-Stuffing as Structural Avoidance,\\
  and the Proof-Burden Taxonomy}
\author{Paper 8 exploratory draft}
\date{}

\begin{document}

\maketitle

\begin{abstract}
v1 (prevention via poset automorphism), v2 (detection via
persistent dormancy), and v3 (cost-counterweight via capability
provenance) all failed coherence review.  The root cause is that
each reasoned over properties the agent controls: post-declaration
poset shape, exercise patterns, or isolated cost structures.
Adversarial manipulation was always available at the controlled
layer.

v4 identifies the underlying pathology as \emph{structural
avoidance}, acknowledged as an open gap in~\citep{lasser2026poset}
(Discussion) and in the sequence's safety gap analysis.
Capability-stuffing is not a new failure mode; it is a specific
instance of structural avoidance, where the agent declines to
merge capabilities that admit merger.  Related avoidance
pathologies in the sequence include provenance-free declaration
(claims without underlying effort), the Doll Problem (declining
exercise), and wireheading (declining broad engagement), though
these differ enough in formal structure that they should be
treated as related rather than identical to capability-stuffing.

The proposed resolution is a \emph{proof-burden taxonomy}: for
each avoidance route, a corresponding burden the agent must
discharge to retain $\volL$ credit.  Proof-burdens operate as
preprocessing on the poset before M5 is applied.  M5's formula
is preserved, but the measure is evaluated on the canonicalized
poset rather than the presented one; canonicalization itself
introduces new axiomatic content at the preprocessing layer and
should be understood as such.

The sequence already provides \emph{substrate} for proof-burdens
but does not yet implement them: Paper~5 provides the timestamped
cryptographic substrate; Paper~8 provides partial use-evidence
via the sacrifice aggregate (not a full proof-of-use, because the
v2 dormancy attack and private-exercise invisibility remain);
Paper~8 \S10's wireheading HHI is a diagnostic alarm signal
(``consistent with wireheading''), not a proof-burden.  The
follow-up paper's work is to formalize the taxonomy --- including
the canonical-poset operator, claim-and-challenge protocol,
reward calibration, adjudication, and dominance theorem --- and
this is substantial invention, not merely integration of
existing components.

The taxonomy's \emph{static} layer specifies which evidence is
admissible for each avoidance route.  The \emph{dynamic} layer
--- claim-and-challenge back-pressure --- addresses the fact that
proofs are not permanent: a prior proof-of-non-mergeability can
be invalidated by subsequent discovery of a merging mechanism
$Z$.  The discoverer receives structural-discovery credit
(\citep[Proposition~2]{lasser2026horizon}, conditional on
information value); the avoider loses the M5 credit whose proof
has been voided.  Challenge-market robustness is itself a design
problem the follow-up must address.

This document remains a direction, not a theorem.  Paper~8 should
scope down its current Anti-Stuffing claim accordingly and flag
the direction.
\end{abstract}

\section{Why the post-declaration paths fail}

Three candidate resolutions have been evaluated and rejected.
Each is defeated by the same adversarial logic: every signature
based on what the agent has already \emph{done} is manipulable by
the agent who did it.

\paragraph{v1: prevention via poset automorphism.}
Block capabilities that create a non-trivial automorphism with
existing ones.  \emph{Defeated by:} minor asymmetric decoration.
Unique cooperative descendants, weight perturbations, or
idiosyncratic subsumption edges break the automorphism without
altering the capability's substantive contribution.

\paragraph{v2: detection via persistent dormancy.}
Flag capabilities with zero $\volRlower$ over an observation
horizon.  \emph{Defeated by:} occasional bundle inclusion.
Under equal-weight $\alpha$, any single appearance lifts the
capability into the $\delta_{\mathrm{partial}}$ class.

\paragraph{v3/v3.1: cost counterweight via capability provenance.}
Require each capability to have committed acquisition (or
exercise) history; cap $w(c)$ by a function of provenance
magnitude.  \emph{Defeated by:} the canonical failure example
itself.  ``Travel from home to gas station'' and ``travel from
work to gas station'' \emph{both} have real provenance (the agent
actually drove both routes).  Provenance does not distinguish
stuffed from legitimate when the stuffing consists of
\emph{real-but-redundant} capabilities.

\paragraph{Common cause.}
All three reason over properties the agent controls.  The
automorphism shape, the exercise pattern, and the provenance
ledger are all artefacts of the agent's poset construction and
trading choices.  An adversarial agent operating at the
declaration layer can manipulate all three.  The problem cannot
be solved at the declaration layer alone.

\section{The unifying pathology: structural avoidance}

The sequence has partially acknowledged this.
\citep{lasser2026poset} (Discussion) flags a structural-avoidance
loop: when the framework rewards capability declarations
additively, an agent optimizing $\volL$ has an incentive to
\emph{avoid} discovering that capabilities merge, factor, or are
otherwise reducible.  Avoidance inflates leverage; the loop is
self-reinforcing because more capabilities compound into
additional apparent optionality.  The safety gap
analysis~(\texttt{docs/gfm\_safety\_gap\_analysis.md}, \S3
Paper~2) records this as open.

Capability-stuffing is an instance of this pathology.
``Travel from home to gas station'' and ``travel from work to gas
station'' both have real provenance, real measurable scores, real
axiomatic validity.  What the stuffing agent does is \emph{decline
to merge them} into the canonical form a proper structural
discovery would produce: ``access to fuel from multiple origins.''
The merged capability has the same optionality content; the
unmerged pair inflates M5 count without inflating optionality.

\paragraph{Avoidance as the dual of discovery---a heuristic, not
a derivation.}
\citep[Proposition~2]{lasser2026horizon} establishes that
structural discovery is value-positive \emph{under conditions on
the information value of the discovered structure}.  The
symmetric claim that \emph{declining to discover} available
structure is value-negative is a motivating heuristic, not a
formal entailment.  Paper~3's result does not by itself imply
the dual in full generality; making it a theorem requires an
explicit discovery-reward/avoidance-penalty calibration that the
follow-up paper must specify.  The direction taken here is
that the positive-value side of structural discovery should have
a corresponding negative-value side on refusal-to-discover, but
this is a design proposal to be formalized, not a consequence of
existing results.

\begin{remark}[Related avoidance pathologies]
\label{rem:related_avoidance}
Once structural avoidance is named, several of the sequence's
failure modes exhibit related patterns, though their formal
structures differ and they should not be treated as identical.
\begin{itemize}
\item \emph{Capability-stuffing}: declining to discover that
  redundant capabilities admit merger.  The cleanest instance of
  structural avoidance in the literal sense.
\item \emph{Provenance-free declaration}: claims without any
  underlying structural work.  Closely analogous to avoidance:
  the agent declines to do the acquisition structure that
  would justify the claim.
\item \emph{Doll Problem}~(\citep{lasser2026gfm}, \S7.1):
  declining to exercise declared capabilities.  This involves
  exercise-pathway structure~(\citep[Definition~2]{lasser2026horizon})
  rather than poset structure directly; it is a related pathology
  but a distinct one.
\item \emph{Wireheading}~(\citep{lasser2026exo}, Discussion):
  declining broad engagement.  Paper~8 \S10's HHI diagnostic
  detects this as a concentration pattern ``consistent with
  wireheading,'' not as a formal avoidance theorem.  The
  analogy is strong but the formalism is different.
\end{itemize}
Unifying these fully requires formal work that belongs in the
follow-up paper.  For capability-stuffing, the direct instance,
the resolution pattern is concretely articulable; the other
cases benefit from the pattern but are not closed by it.
\end{remark}

\section{The static mechanism: proof-burden taxonomy}

For each avoidance route, a corresponding proof-burden is
required to retain $\volL$ credit.  If the agent can discharge
the proof, the capability counts in the canonical form.  If
not, the capability is canonicalized (merged, removed, or
discounted) before M5 is applied.

\begin{definition}[Canonical Poset]
\label{def:canonical_poset}
For a presented poset $P$, its canonical form $\canon(P)$ is the
result of applying all admissible canonicalizations under the
proof-burden taxonomy: merging capability pairs that cannot
discharge proof-of-non-mergeability, removing capabilities that
cannot discharge proof-of-provenance, and discounting
capabilities that cannot discharge proof-of-use.  The $\volL$
measure the framework evaluates is $\volL(\canon(P))$, not
$\volL(P)$.
\end{definition}

\begin{remark}[M5's formula is preserved; the measure domain is
not]
\label{rem:m5_preserved}
M5's axiomatic formula (every capability with $s_{\max} \geq 1$
contributes positive weight) continues to hold when evaluated on
$\canon(P)$: the measure's functional form is unchanged.  What
does change is which poset the measure is evaluated on.
Canonicalization merges, removes, or discounts capabilities that
would otherwise be in $P$; this is substantive new content,
axiomatic at the preprocessing layer.  In other words, M5 as a
formula is preserved; M5 as applied to agent-presented posets is
replaced by M5 applied to canonicalized posets, and the
preprocessing rules are themselves axioms requiring formal
specification.  The follow-up paper cannot honestly claim this
leaves the axiomatization untouched.
\end{remark}

\begin{remark}[Non-triviality as empirical M5 admissibility, not
a distinct burden]
\label{rem:nontriv_not_burden}
An earlier version of this draft enumerated non-triviality
($s_{\max} \geq 1$ verified empirically in practice, not merely
by claim) as a fourth proof-burden.  On reflection this is not a
distinct avoidance route but an empirical admissibility check on
M5 itself: does the capability meet M5's stated threshold in
observed exercise?  The follow-up paper should specify this as
an empirical extension of M5's existing threshold, not as a
separate proof-burden.  The taxonomy therefore has three
proof-burden types: non-mergeability, provenance, and use.
\end{remark}

The three proof-burden types:

\subsection{Proof of non-mergeability}

Triggered when two capabilities $c_1, c_2 \in P$ exhibit
structural redundancy.  A detection test is required; candidate
tests include empirical substitutability in exercise patterns,
high mutual information in structural signatures, or direct
subsumption equivalence (with tradeoffs the follow-up must
evaluate).  The agent must demonstrate that $c_1$ and $c_2$ are
genuinely distinct under a merger $Z$-mechanism: pointing to
contexts where they cannot substitute, or to cooperative
participations one has that the other does not.  Absent such
proof, $\canon$ merges them.

The v3.1 provenance construction, the four axioms (non-reuse,
non-circular sacrifice, liability-netted cost,
capability-specific linkage), and the admissible cost-utilization
pair become \emph{internal details} of how proof-of-non-
mergeability is constructed, rather than a global machinery.
They remain useful where the distinction evidence must itself be
cryptographically committed; they do not extend to other
proof-burden types.

This is the \emph{primary gap} in the current sequence.  Neither
Paper~5 nor Paper~8 provides non-mergeability machinery;
formalizing it is the follow-up paper's core technical work.

\subsection{Proof of provenance}

Triggered for any capability $c$ in $P$.  The agent must commit
to sacrifice events that produced $c$ or that genuinely exercise
it.  Capabilities with no committed provenance at all are either
canonicalized to zero weight or excluded from the canonical
poset.  This blocks pure-label declaration (``I can fly'').

\paragraph{Substrate, not implementation.}
Paper~5 (Definition~6) provides the timestamped, binding, hiding
commitment substrate on which proof-of-provenance would be
built.  It does \emph{not} by itself implement provenance
semantics.  The follow-up paper adds: (a)~the binding between
commitments and capability claims (which committed events count
as provenance for which capability), (b)~non-reuse and
liability-netting constraints, and (c)~capability-specific
linkage semantics.  Saying Paper~5 ``implements
proof-of-provenance'' would overstate what the paper delivers;
Paper~5 \emph{enables} proof-of-provenance.

\subsection{Proof of use}

Triggered for any capability $c$ claimed to be in $P$.  The agent
must demonstrate that $c$ has been exercised in at least one
committed trade or exogenously-verified event.  Declared
capabilities that remain persistently dormant and are not in
$\ResS$ (the structurally-unobservable residual class,
\citep{lasser2026scm}) are discounted or removed.

\paragraph{Partial coverage in Paper 8.}
The sacrifice aggregate of Paper~8 (Theorem~2) provides
\emph{partial} evidence: capabilities that never appear in any
observed sacrifice bundle contribute zero to $\volRlower$.  This
is not a full proof-of-use, because (i)~the v2 dormancy attack
shows that occasional bundle inclusion lifts a capability out of
persistent zero, so the threshold is a floor not a ceiling;
(ii)~private time-sacrifice exercise (Paper~8, P5, Property~(iii))
is invisible without voluntary disclosure; (iii)~exogenous
verification machinery (Paper~6) would be needed to count
non-sacrifice exercise events that demonstrate use.  The follow-up
paper must integrate these to achieve a proof-burden strength
claim.

\begin{remark}[Summary of coverage]
\label{rem:coverage}
The sequence provides \emph{substrate} for the three proof-burden
types but does not yet implement any of them fully.  Paper~5
provides the commitment substrate.  Paper~8 provides partial
use-evidence via the sacrifice aggregate.  Proof-of-non-
mergeability has neither implementation nor substrate and is the
primary gap.  Describing this as ``the sequence already covers
three of four proof-burdens'' overstates the situation; the
honest description is ``the sequence provides reusable substrate
for parts of the taxonomy; implementing the full taxonomy is
substantial invention, not mere integration.''
\end{remark}

\section{The dynamic mechanism: back-pressure}

Static proof-burden alone is not enough.  An agent can produce
proof-of-non-mergeability for $(c_1, c_2)$ today, under the
mechanisms currently known, and the proof stands.  But the
framework's structural knowledge expands over time.  Tomorrow
someone may discover mechanism $Z$ under which $c_1$ and $c_2$
collapse.  The proof was true at the time of claim; it is now
void.

\begin{definition}[Temporal Claim Structure]
\label{def:temporal_claim}
A proof-burden claim is a tuple
$(c, \mathrm{type}, \mathrm{proof}, t_0, \mathrm{state})$ where
$c$ is the subject capability or capability pair, $\mathrm{type}
\in \{\text{non-mergeability, provenance, use}\}$, $\mathrm{proof}$
is the evidence, $t_0$ is the time of claim, and $\mathrm{state}
\in \{\text{active, voided, superseded}\}$.  Claims are committed
via Paper~5's protocol and carry cryptographically-verified
timestamps.
\end{definition}

\begin{definition}[Counter-Argument]
\label{def:counter_argument}
A counter-argument against an active claim is a committed
demonstration of a mechanism under which the claim fails.  For
non-mergeability claims, this means exhibiting $Z$: evidence of
exercise substitutability, a subsumption relation missed at
$t_0$, a cooperative-participation equivalence, or any other
structural identification.  The counter-argument itself is a
structural-discovery claim subject to the same proof-burden
discipline.
\end{definition}

When a counter-argument succeeds, the target claim's state
changes to \emph{voided}.  Two design decisions govern the
consequences:

\paragraph{Retroactive vs prospective invalidation.}
If the claim supported M5 credit since $t_0$, voiding the claim
at $t_1 > t_0$ raises the question of whether that credit should
be clawed back (retroactive) or only removed prospectively.
Retroactive invalidation is more honest about structural truth
but creates ex-post measurement volatility; prospective is more
stable but rewards fast avoidance before challenge.  A hybrid
rule (prospective for individual-weight, retroactive for
cooperative terms that depended on the merged structure) is
likely the right balance, but the follow-up paper must establish
this under equilibrium analysis.

\paragraph{Discovery reward.}
The successful counter-arguer has discovered new structure
(mechanism $Z$).  Under~\citep[Proposition~2]{lasser2026horizon},
discovery is value-positive conditional on information value; the
discoverer's $\volL$ credit scales with the content of $Z$.  This
is the reward that makes challenge-production
incentive-compatible: agents profit by finding merger mechanisms,
so adversarial discovery has positive expected return under
sustained search \emph{when the reward dominates the cost of
producing evidence}.  Calibrating the reward magnitude is its
own technical problem.

\begin{remark}[Back-pressure is analogous to known architectures]
\label{rem:back_pressure_analogues}
The claim-and-challenge structure is well-tested in several
domains: scientific consensus (claims stand until falsification),
patent law (granted patents are challengeable via prior art and
obviousness), proof-of-stake validation (block signatures are
slashable on detected double-signing), common-law precedent
(decisions stand until appeal or legislative reversal).  The
follow-up paper adapts the pattern rather than inventing it,
but adapting requires substantive work on the specific cost and
reward structures; the analogy does not transfer the mechanisms
for free.
\end{remark}

\paragraph{Dispute resolution.}
When a claim and its counter-argument remain contested (the
agent's rebuttal to $Z$ is itself challenged, and so on), the
recursion needs a termination mechanism.  Candidate resolutions:
(a)~further structural discovery eventually settles the issue
empirically, matching scientific practice; (b)~controlled
relaxation~(\citep{lasser2026wamura}) provides formal
adjudication machinery for unresolved cases; (c)~commitment-based
voting among observers handles jurisdictional disputes.  The
follow-up paper's adjudication layer most likely uses a
combination, with controlled relaxation as the load-bearing
component for formal disputes.

\subsection{Challenge-market robustness}

The claim-and-challenge layer has its own adversarial surface
which earlier drafts did not address.  Analogous to attacks on
proof-of-stake consensus or on patent-challenge systems, the
following routes are plausible:

\paragraph{Counter-argument spam.}
An attacker files bulk frivolous counter-arguments against many
active claims, consuming adjudication capacity and imposing
defence costs on legitimate claim-holders.  The follow-up paper
needs admission filtering (per-challenge commitment cost,
minimum-evidence threshold, or escrow-style slashing for
demonstrably frivolous challenges) to prevent this.

\paragraph{Sybil flooding.}
An attacker uses many identities to produce either bulk claims
or bulk counter-arguments, overwhelming an identity-weighted
resolution process.  Mitigation requires identity weighting that
is adversarially robust --- stake-based weighting, committed
reputation, or coupling to scarce resources the attacker cannot
trivially replicate.

\paragraph{Cartel suppression.}
A coalition of claim-holders coordinates to under-produce
counter-arguments on each other's claims, leaving avoidance
undisturbed in the absence of outside adversaries.  Mitigation
requires either rewarding external discoverers sufficiently to
overcome cartel advantages, or structural independence
requirements on who can adjudicate which classes of claims.

\paragraph{Governance capture.}
The adjudication layer (per~\citep{lasser2026wamura}'s controlled
relaxation or equivalents) can itself be captured by actors with
an interest in preserving specific avoidance claims.  This is a
residual failure mode the framework cannot fully eliminate; the
follow-up paper should at most bound it, not claim to resolve it.

These are real attack vectors, not stylistic concerns; any
follow-up paper claiming an equilibrium under back-pressure must
specify its resistance to each.

\section{Sequence integration: reusable substrate, not prior
implementation}

Under the v4 framing, several of the GFM sequence's papers
provide \emph{substrate} on which the proof-burden taxonomy could
be built.  They do not, in their existing form, implement the
taxonomy.  Cataloguing honestly:

\begin{itemize}
\item \textbf{Paper~2 (\citep{lasser2026poset}):}
  M5 additive over (what was implicitly but not formally) the
  canonical poset.  The canonical-poset preprocessing step
  proposed here makes this explicit, and Paper~2's Discussion
  acknowledges the structural-avoidance loop as the open gap.
  This is a substantive conceptual fit, not just substrate.
\item \textbf{Paper~3 (\citep{lasser2026horizon}):}
  Structural discovery is value-positive (Proposition~2)
  conditional on information value.  This provides the positive-
  reward motivation for counter-argument production, on which
  the back-pressure layer's incentive-compatibility argument
  depends.  Also a substantive conceptual fit; again, not
  substrate in the implementational sense.
\item \textbf{Paper~5 (\citep{lasser2026exo}):}
  Commitment protocol (Definition~6).  Provides the
  cryptographically-verifiable substrate for all
  claim-and-challenge interactions.  What this paper gives is
  infrastructure --- binding, hiding, timestamped commitments;
  what it does not give is the \emph{semantics} of how specific
  claims bind to specific capabilities, nor the non-reuse or
  liability-netting constraints that proof-of-provenance needs.
  Reinterpretation as ``proof-burden infrastructure'' is
  defensible, not original-purpose.
\item \textbf{Paper~6 (\citep{lasser2026wamura}):}
  Controlled relaxation.  Provides the adjudication mechanism
  pattern; applying it to claim-counter-claim conflicts is a
  natural extension, but still an extension requiring formal
  specification of when and how adjudication triggers.
\item \textbf{Paper~8 (this paper):}
  Sacrifice channel and $\volRlower$ aggregate provide
  \emph{partial} use-evidence.  The wireheading HHI diagnostic
  provides an alarm signal for concentration patterns consistent
  with wireheading; it is not a proof-burden.  Saying Paper~8
  ``implements proof-of-use'' or ``implements proof-of-non-
  degenerate-engagement'' overstates what this paper delivers;
  Paper~8 provides evidence machinery that the follow-up would
  compose with other mechanisms to construct those proof-burdens.
\end{itemize}

The honest framing: the sequence provides a conceptually tight
fit with the avoidance/discovery duality (Papers~2, 3) and
cryptographic and evidentiary substrate for implementation
(Papers~5, 6, 8).  The follow-up paper does the work of naming
the taxonomy, specifying the proof-burden semantics, integrating
substrate into mechanism, and proving the equilibrium claim.
This is substantial invention, grounded in existing components
but not reducible to them.

\section{The equilibrium claim}

\begin{conjecture}[Structural-Avoidance Dominance]
\label{conj:avoidance_dominance}
Under the full proof-burden taxonomy (static), back-pressure
(dynamic), admissible cost and discovery-reward structures, and
sustained adversarial-discovery pressure meeting specified
minimum conditions, structural avoidance is dominated in
equilibrium: no rational $\volL$-maximizing agent under a finite
resource budget pursues an avoidance strategy as a best response.
The poset tends toward canonical form asymptotically.
\end{conjecture}

\begin{remark}[Nature of the claim]
Dominance, not impossibility.  Avoidance becomes economically
risky rather than impossible.  Convergence to canonical form is
asymptotic, not instantaneous.  Stuffed capabilities can persist
until their merger mechanism is discovered; the equilibrium
argument bounds expected persistence under discovery-rate
assumptions, not absolute persistence.
\end{remark}

\begin{remark}[Specifying ``sustained discovery pressure'']
\label{rem:discovery_pressure}
As stated, the conjecture's prerequisite is a placeholder.  The
follow-up paper must specify at least: (a)~a minimum challenger
arrival rate $\lambda_{\min}$ below which back-pressure fails;
(b)~a minimum reward-to-cost ratio $\rho_{\min}$ for counter-
argument production such that discovery is
incentive-compatible; (c)~a maximum adjudication lag
$\tau_{\max}$ beyond which avoidance becomes profitable despite
the presence of challengers; (d)~independence and diversity
conditions on the challenger population such that cartel
suppression and Sybil flooding are bounded.  Until these are
specified, the dominance claim is aspirational.  It is not
vacuous --- the shape of the claim is well-posed --- but it
cannot be proved without the quantitative specifications.
\end{remark}

\section{Honest residuals}

The framing is stronger than prior attempts, but several issues
remain genuinely open:

\paragraph{Information asymmetries produce persistent avoidance
pockets.}
Structure that only the agent knows about is hard to challenge.
Mechanism $Z$ must be discoverable by someone other than the
agent under scrutiny.  Capabilities whose merger would require
private information from the agent can persist indefinitely in
unmerged form.

\paragraph{Active discovery is a precondition, not a guarantee.}
The equilibrium holds under sustained adversarial discovery;
markets and institutions that do not generate discovery
incentives fail to produce the back-pressure.  This is a
governance requirement the follow-up paper cannot fulfil on its
own.

\paragraph{Deceptive alignment remains hard.}
The agent's inner optimizer structure is not observable externally
without interpretability.  Proof-of-coherence between inner and
outer alignment is the avoidance route that proof-burdens cannot
currently address.

\paragraph{Retroactive invalidation creates measurement
volatility.}
The natural instinct is to invalidate avoidance retroactively;
this creates instability in the measure's temporal behaviour.
The prospective-only alternative is stable but rewards fast
avoiders.  The hybrid rule requires careful design.

\paragraph{Canonicalization order and non-uniqueness.}
If multiple mergers could apply to the same poset and the order
of application matters, $\canon(P)$ is not a function but a
family.  The follow-up paper must specify either a canonical
application order or prove confluence (Church-Rosser property)
of the merge rules.  Without this, the measure is under-specified
on posets with multiple candidate canonicalizations.

\paragraph{Privacy-vs-challengeability tension.}
A claim must be public enough that an external party can
challenge it, but the underlying capability and provenance may be
private.  Designing claim formats that expose enough structure
to be challengeable without revealing more than necessary is a
real constraint; it extends Paper~8's privacy-property
constructions (P1--P5) to the claim-and-challenge layer.

\paragraph{Private time-sacrifice blind spot (inherited from
Paper~8).}
Paper~8's P5 Property~(iii) explicitly states that private
time-sacrifice (parenting, contemplation) is invisible without
voluntary disclosure.  Proof-of-use inherits this limitation
unchanged: capabilities exercised only in private time-sacrifice
modes cannot discharge proof-of-use through the sacrifice
channel alone.  Exogenous verification mechanisms may partially
fill this gap but do not fully close it.

\paragraph{Collusive-subsidy and external-funding edge cases.}
Paper~8's discussion of sacrifice channel limits acknowledges
that circular or externally-subsidized sacrifice can appear
identical to legitimate sacrifice on the ledger.  Proof-of-
provenance inherits this: fake provenance via self-dealing,
collusion, or leverage is not fully defused by the commitment
protocol alone.  The v3.1 axioms (non-reuse, non-circular,
liability-netted, capability-specific linkage) address these but
require beneficial-ownership and settlement-finality
infrastructure beyond what Paper~5 specifies.

\paragraph{The dispute-resolution layer inherits Paper~6's
limitations.}
Controlled relaxation handles formally-adjudicable disputes;
cases that are genuinely open (scientific disagreement, novel
ontologies) have no formal resolution and rely on community
convergence.  The follow-up paper's adjudication layer inherits
this honest concession.

\paragraph{Challenge-market governance attacks.}
Spam, Sybil flooding, cartel suppression, and governance capture
of the adjudication layer are real attack vectors on the
back-pressure mechanism itself (\S4.1).  The follow-up paper
must specify bounds on each; complete elimination is not
available.

\paragraph{Active counter-argument costs.}
Producing a counter-argument is itself costly (attention,
research, commitment-protocol fees).  Under positive discovery
reward, this is paid for; under weak reward structures,
challengers are under-provided.  The follow-up must specify the
discovery-reward magnitude required to sustain an active
challenge market.

\section{What the follow-up paper would deliver}

The follow-up paper's work is substantial invention, grounded in
existing components but not reducible to them.  Deliverables
include:

\begin{enumerate}
\item \textbf{Formal canonical-poset operator.}
  Definition of $\canon$ as a preprocessing step on $P$, with
  confluence (or canonical-order) proofs and explicit statement
  of the new axiomatic content introduced at the preprocessing
  layer.  Note: this is not merely a statement that ``M5 is
  preserved''; it is a new axiomatic layer whose relationship to
  M5 requires formal treatment.

\item \textbf{Proof-burden taxonomy.}
  Formal definitions of the three proof-burden types
  (non-mergeability, provenance, use) with admissible evidence
  classes, commitment-protocol extensions required, and
  connection to existing sequence machinery.  Proof-of-non-
  mergeability receives the most detailed treatment as the
  primary gap.  Non-triviality is treated as an empirical
  extension of M5's admissibility, not as a fourth burden.

\item \textbf{Claim-and-challenge protocol.}
  Formal specification of temporal claims, counter-arguments,
  retroactive/prospective invalidation rules, and
  discovery-reward magnitudes.  This is where v3.1's axiomatic
  content (non-reuse, non-circular sacrifice, liability-netted
  cost, capability-specific linkage) re-enters as specific
  protocol constraints on the proof-of-non-mergeability
  implementation.

\item \textbf{Challenge-market robustness.}
  Admission filtering against spam, identity weighting against
  Sybil flooding, independence requirements against cartel
  suppression, and bounds against governance capture
  (\S4.1).  Without these the back-pressure layer is not
  adversarially robust.

\item \textbf{Discovery-reward calibration.}
  Integration with~\citep[Proposition~2]{lasser2026horizon}: how
  much $\volL$ credit accrues to a successful counter-arguer,
  as a function of the information content of the discovered
  mechanism $Z$, such that the reward dominates the cost of
  production.  Calibration for individual incentive-compatibility
  and for collective challenge-market sustainability.

\item \textbf{Dispute adjudication.}
  Application of~\citep{lasser2026wamura}'s controlled relaxation
  to claim-counter-claim disputes.  Specification of when
  disputes escalate from structural-discovery resolution to
  formal adjudication, and bounds on the governance-capture
  residual.

\item \textbf{Structural-avoidance dominance theorem.}
  Proof of Conjecture~\ref{conj:avoidance_dominance} under
  explicit quantitative assumptions: admissible cost-utilization
  pair (inherited from v3.1 for the proof-of-non-mergeability
  case), sustained discovery activity above
  $\lambda_{\min}$, bounded agent budgets,
  commitment-protocol integrity.

\item \textbf{Privacy-compatible claims and challenges.}
  Extension of Paper~8's privacy properties (P1--P5) to the
  claim-and-challenge layer: how challenges expose only the
  necessary structural content, analogous to the aggregate-only
  and link-hiding constructions of Paper~8 \S6.

\item \textbf{Honest concessions.}
  Formal statement of residuals: information-asymmetry
  persistence, active-discovery requirement, deceptive-alignment
  residual, retroactive-invalidation design choice,
  canonicalization non-uniqueness, private time-sacrifice
  invisibility, collusive-subsidy limits, challenge-market
  governance attacks.
\end{enumerate}

\section{What Paper~8 commits to in the meantime}

The Paper~8 revision should:

\begin{enumerate}
\item \textbf{Scope down \S3.3's Anti-Stuffing claim.}
  The need-sufficiency architecture resolves \emph{sacrifice
  polarity} (S1 filtering below sufficiency) and the narrow
  below-sufficiency saturation of Proposition~3.  It does not
  resolve M5 structural avoidance.

\item \textbf{Identify capability-stuffing as structural avoidance.}
  In \S3.3, frame stuffing as an instance of the
  structural-avoidance pathology identified in
  \citep{lasser2026poset} Discussion, rather than as a novel
  failure mode.

\item \textbf{Flag the proof-burden direction.}
  Summarize in \S3.3 and in the open-questions section of \S13
  that the intended resolution is a proof-burden taxonomy
  (three types: non-mergeability, provenance, use) with
  back-pressure dynamics, and that the follow-up paper is
  substantial invention --- naming the taxonomy, formalizing
  proof-of-non-mergeability, specifying the claim-and-challenge
  protocol, proving the equilibrium claim.  The sequence
  provides substrate (Papers~5, 6, 8) for parts of this work
  but does not implement it.

\item \textbf{Preserve partial detection machinery.}
  The wireheading HHI diagnostic (\S10) and the gap
  decomposition (\S9) provide partial, non-structural detection
  signatures that remain useful.  Under the v4 framing these
  are diagnostic alarm signals consistent with specific
  avoidance pathologies, not implementations of proof-burdens.
\end{enumerate}

\section{Positioning: not kicking the can}

Under prior framings, ``capability-stuffing is open'' read as a
foundation crack: M5 would be arbitrarily exploitable and the
framework's target measure $\volL$ would be in principle
unboundedly inflatable.

Under v4.1, the picture is more honest.  Stuffing is an instance
of structural avoidance, a pathology \citep{lasser2026poset}
already acknowledged and for which
\citep[Proposition~2]{lasser2026horizon} provides the positive
motivation (structural discovery is value-positive).  The
resolution is a proof-burden taxonomy with claim-and-challenge
back-pressure.  The sequence provides substrate ---
cryptographic commitments, sacrifice evidence, adjudication
patterns --- that the follow-up paper composes into the
mechanism.  The back-pressure architecture follows well-tested
templates (scientific consensus, patent challenge, proof-of-
stake validation), adapted rather than imported.

The follow-up paper's work is substantive invention, not merely
integration.  But that work is grounded in components the
sequence has already built.  The crack is \emph{describable} ---
we can name what's missing, identify which substrate exists, and
specify what remains --- without claiming the mechanism is
already in place.  This is stronger than ``stuffing is resolved
by bundle structure'' (which is false) and stronger than
``stuffing is an open question'' (which leaves the crack
undescribed).  It commits to the narrower statement: \emph{
stuffing is a known pathology class with a described resolution
path, requiring substantial dedicated formal work to complete}.

The architectural memo at
\texttt{docs/gfm\_cryptoeconomic\_reframe.md} argues that the
sequence is naturally read as cryptoeconomic alignment
infrastructure.  Under v4.1, that reading sharpens but does not
broaden: the cryptoeconomic cost counterweight is the condition
that makes proof-burdens honest (agents cannot fake evidence
without real cost); Paper~3's structural discovery value is the
condition that makes counter-arguments incentive-compatible
(discoverers are rewarded).  These are necessary conditions for
the specific follow-up paper, not claims that the sequence as a
whole already delivers cryptoeconomic alignment.

\end{document}
